% =====================================================================
\chapter{Soluciones de los ejercicios}
\label{ap:soluciones}
% =====================================================================

Cada solución se presenta en cinco partes: datos, planteamiento, desarrollo, resultado y comprobación. En los pasos intermedios se redondea a cuatro cifras significativas salvo que se indique otra cosa; los resultados exactos se dan como fracción o con radicales cuando existen.

% ---------------------------------------------------------------------
\section{Tema 1}

\begin{solucion}{eje:1-vf}
\planteamiento contrastar cada afirmación con los resultados del \cref{tema:1}.
\pasos
\begin{enumerate}[label=(\alph*)]
  \item \textbf{Verdadera} si $f$ es diferenciable. Si $r(t)$ recorre una curva de nivel, $f(r(t))=c$ y el \cref{cor:nivel} da $\nabla f(r(t))\T r'(t)=0$: el gradiente es ortogonal a la tangente.
  \item \textbf{Verdadera.} Por el \cref{teo:max-ascenso}, $D_uf\ge-\norm{\nabla f}$, con igualdad solo en $u=-\nabla f/\norm{\nabla f}$.
  \item \textbf{Falsa.} La función $f(x,y)=xy/(x^{2}+y^{2})$, con $f(0,0)=0$, tiene ambas parciales nulas en el origen, pero $f(t,t)=1/2$ para $t\ne 0$; no es continua en $(0,0)$ y, por tanto, no es diferenciable allí.
\end{enumerate}
\resultado \textbf{(a) verdadera; (b) verdadera; (c) falsa.}
\comprobacion en (c), si $f$ fuera diferenciable en el origen, su gradiente sería $(0,0)\T$ (\cref{prop:grad-parciales}) y se tendría $f(t,t)=o(\abs{t})\to 0$, en contradicción con $f(t,t)=1/2$.
\end{solucion}

\begin{solucion}{eje:1-gradiente}
\datos $f(x,y)=x^{3}+2xy^{2}$; punto $(1,2)$; $u=(1,1)\T/\sqrt2$.
\planteamiento derivar respecto de cada variable con la otra fija y aplicar la \cref{prop:direccional}.
\pasos
\begin{enumerate}[label=\arabic*.]
  \item $\partial f/\partial x=3x^{2}+2y^{2}$; en $(1,2)$: $3+8=11$.
  \item $\partial f/\partial y=4xy$; en $(1,2)$: $8$.
  \item $\nabla f(1,2)=(11,8)\T$ y $\norm{\nabla f(1,2)}=\sqrt{121+64}=\sqrt{185}\approx 13{,}60$.
  \item $D_uf(1,2)=(11+8)/\sqrt2=19/\sqrt2\approx 13{,}44$.
\end{enumerate}
\resultado \textbf{$\nabla f(1,2)=(11,8)\T$; $\norm{\nabla f}=\sqrt{185}\approx 13{,}60$; $D_uf=19/\sqrt2\approx 13{,}44$.}
\comprobacion $13{,}44<13{,}60$, como exige el \cref{teo:max-ascenso}. Las dos cifras son próximas porque el ángulo entre $u$ y $\nabla f$ es pequeño: su coseno es $13{,}44/13{,}60\approx 0{,}988$.
\end{solucion}

\begin{solucion}{eje:1-d2l}
\datos $f(x)=3x_1^{2}+5e^{x_2}$; punto $(1,0)$; dirección $(3,4)\T$.
\planteamiento calcular el gradiente, normalizar la dirección y usar $D_uf=\nabla f\T u$.
\pasos
\begin{enumerate}[label=\arabic*.]
  \item $\nabla f(x)=(6x_1,\;5e^{x_2})\T$.
  \item En $(1,0)$: $\nabla f=(6,\;5e^{0})\T=(6,5)\T$.
  \item $\norm{(3,4)}=5$, luego $u=(0{,}6;\,0{,}8)\T$.
  \item $D_uf(1,0)=6\cdot 0{,}6+5\cdot 0{,}8=3{,}6+4=7{,}6$.
  \item $\norm{\nabla f(1,0)}=\sqrt{36+25}=\sqrt{61}\approx 7{,}810$.
\end{enumerate}
\resultado \textbf{$\nabla f(1,0)=(6,5)\T$; $D_uf=7{,}6\le\sqrt{61}\approx 7{,}810$.}
\comprobacion la diferencia centrada en $x_2$ con $h=0{,}01$ da $5\,(e^{0{,}01}-e^{-0{,}01})/0{,}02=5\cdot 1{,}0000167\approx 5{,}0001$.
\end{solucion}

\begin{solucion}{eje:1-norma}
\datos $f(x)=\norm{x}=\bigl(\sum_ix_i^{2}\bigr)^{1/2}$; $x=(3,4)\T$.
\planteamiento regla de la cadena en una variable con $s=\sum_ix_i^{2}$ y $f=s^{1/2}$.
\pasos
\begin{enumerate}[label=\arabic*.]
  \item $\partial s/\partial x_i=2x_i$ y $d(s^{1/2})/ds=\tfrac12 s^{-1/2}$.
  \item $\partial f/\partial x_i=\tfrac12 s^{-1/2}\cdot 2x_i=x_i/\norm{x}$, luego $\nabla f(x)=x/\norm{x}$, el vector unitario en la dirección de $x$. Las parciales son continuas fuera del origen, luego $f$ es diferenciable allí (\cref{teo:C1}).
  \item En $(3,4)$: $\norm{x}=5$ y $\nabla f=(0{,}6;\,0{,}8)\T$.
  \item En $x=0$: $f(t\,e_1)=\abs{t}$, cuyo cociente $\abs{t}/t$ vale $-1$ para $t<0$ y $+1$ para $t>0$. La parcial $\partial f/\partial x_1(0)$ no existe y $f$ no es diferenciable en el origen.
  \item El recorte por norma sustituye $g$ por $c\,g/\norm{g}$ cuando $\norm{g}>c$: conserva la dirección $\nabla\norm{g}=g/\norm{g}$ y fija la longitud en $c$.
\end{enumerate}
\resultado \textbf{$\nabla\norm{x}=x/\norm{x}$ para $x\ne 0$; en $(3,4)$ vale $(0{,}6;\,0{,}8)\T$; en $x=0$ no existe.}
\comprobacion $\norm{\nabla f(x)}=1$ para todo $x\ne 0$: alejarse radialmente del origen aumenta la norma a razón de una unidad por unidad recorrida.
\end{solucion}

\begin{solucion}{eje:1-mit}
\datos $f(x,y)=x^{3}-3x^{2}+xy-y^{2}+x+y$; punto $(1,1/2)$.
\planteamiento por el \cref{teo:max-ascenso}, el máximo se alcanza en $\nabla f/\norm{\nabla f}$, el mínimo en la dirección opuesta, y la derivada direccional es nula en las direcciones ortogonales al gradiente.
\pasos
\begin{enumerate}[label=\arabic*.]
  \item $f_x=3x^{2}-6x+y+1$; en $(1,1/2)$: $3-6+0{,}5+1=-1{,}5$.
  \item $f_y=x-2y+1$; en $(1,1/2)$: $1-1+1=1$.
  \item $\nabla f=(-3/2,\,1)\T$ y $\norm{\nabla f}=\sqrt{9/4+1}=\sqrt{13}/2\approx 1{,}803$.
  \item Máximo: $u=(-3,2)\T/\sqrt{13}\approx(-0{,}8321;\,0{,}5547)$, con $D_uf=\sqrt{13}/2\approx 1{,}803$.
  \item Mínimo: $u=(3,-2)\T/\sqrt{13}$, con $D_uf\approx-1{,}803$.
  \item Nulas: los vectores unitarios ortogonales a $(-3,2)$, es decir,
  \[u=\pm(2,3)\T/\sqrt{13}\approx\pm(0{,}5547;\,0{,}8321).\]
\end{enumerate}
\resultado \textbf{máximo $\sqrt{13}/2\approx 1{,}803$ en $(-3,2)\T/\sqrt{13}$; mínimo $-1{,}803$ en $(3,-2)\T/\sqrt{13}$; nula en $\pm(2,3)\T/\sqrt{13}$.}
\comprobacion $(-3,2)\cdot(2,3)=-6+6=0$, y $\nabla f\T(-3,2)\T/\sqrt{13}=(4{,}5+2)/\sqrt{13}=6{,}5/3{,}606\approx 1{,}803$.
\end{solucion}

\begin{solucion}{eje:1-minimo}
\datos $g=\nabla f(x)\ne 0$; $u$ unitario; $f(x,y)=x^{2}+y^{2}$ y $r(t)=\sqrt c\,(\cos t,\sin t)\T$.
\planteamiento Cauchy--Schwarz para la cota y el \cref{lem:curva} para la ortogonalidad.
\pasos
\begin{enumerate}[label=\arabic*.]
  \item Por la \cref{prop:direccional}, $D_uf(x)=g\T u$.
  \item Por Cauchy--Schwarz, $\abs{g\T u}\le\norm{g}\,\norm{u}=\norm{g}$; en particular, $g\T u\ge-\norm{g}$.
  \item La cota se alcanza: para $u^{-}=-g/\norm{g}$, que es unitario porque $g\ne 0$, se tiene $g\T u^{-}=-\norm{g}^{2}/\norm{g}=-\norm{g}$.
  \item Unicidad: la igualdad en Cauchy--Schwarz exige $u=\alpha g$; con $\norm u=1$, $\alpha=\pm1/\norm{g}$, y solo el signo negativo da $-\norm g$.
  \item Circunferencias: $f(r(t))=c(\cos^{2}t+\sin^{2}t)=c$ para todo $t$, y $r'(t)=\sqrt c\,(-\sin t,\cos t)\T$.
  \item $\nabla f(r(t))=2r(t)=2\sqrt c\,(\cos t,\sin t)\T$, luego $\nabla f(r(t))\T r'(t)=2c(-\cos t\sin t+\sin t\cos t)=0$.
\end{enumerate}
\resultado \textbf{$\min_{\norm u=1}D_uf(x)=-\norm{\nabla f(x)}$, alcanzado solo en $u=-\nabla f/\norm{\nabla f}$; en las circunferencias, el gradiente es radial y ortogonal a la tangente.} $\blacksquare$
\comprobacion la derivada de la constante $f(r(t))=c$ es nula, como debe ser por el \cref{cor:nivel}.
\end{solucion}

% ---------------------------------------------------------------------
\section{Tema 2}

\begin{solucion}{eje:2-formas}
\datos $g\colon\R^{5}\to\R^{3}$, $h\colon\R^{3}\to\R^{2}$, $f\colon\R^{2}\to\R$.
\planteamiento el jacobiano de una función $\R^{n}\to\R^{m}$ es $m\times n$ (\cref{def:jacobiano}).
\pasos
\begin{enumerate}[label=\arabic*.]
  \item $J_g$ es $3\times 5$, $J_h$ es $2\times 3$ y $J_f$ es $1\times 2$.
  \item Por el \cref{teo:cadena}, $J_{f\circ h\circ g}=J_f\,J_h\,J_g$: $(1\times 2)(2\times 3)(3\times 5)=1\times 5$.
  \item El gradiente es la traspuesta: un vector de $\R^{5}$, de forma $5\times 1$.
\end{enumerate}
\resultado \textbf{$3\times 5$, $2\times 3$, $1\times 2$; la composición tiene jacobiano $1\times 5$ y gradiente en $\R^{5}$.}
\comprobacion en cada producto coinciden las dimensiones interiores ($2$ y $3$), y el gradiente tiene la forma de la variable de entrada.
\end{solucion}

\begin{solucion}{eje:2-producto}
\datos $F=uv$ con $u=x+y^{2}$, $v=x^{2}-y$; punto $(1,2)$.
\planteamiento \cref{cor:cadena-escalar}: $\partial F/\partial x=F_u\,u_x+F_v\,v_x$ y análogamente para $y$.
\pasos
\begin{enumerate}[label=\arabic*.]
  \item Valores: $u=1+4=5$ y $v=1-2=-1$.
  \item Derivadas exteriores: $F_u=v=-1$ y $F_v=u=5$.
  \item Derivadas interiores: $u_x=1$, $u_y=2y=4$, $v_x=2x=2$, $v_y=-1$.
  \item $\partial F/\partial x=(-1)\cdot 1+5\cdot 2=9$.
  \item $\partial F/\partial y=(-1)\cdot 4+5\cdot(-1)=-9$.
\end{enumerate}
\resultado \textbf{$\partial F/\partial x=9$ y $\partial F/\partial y=-9$ en $(1,2)$.}
\comprobacion $F=x^{3}-xy+x^{2}y^{2}-y^{3}$, luego $F_x=3x^{2}-y+2xy^{2}=3-2+8=9$ y $F_y=-x+2x^{2}y-3y^{2}=-1+4-12=-9$.
\end{solucion}

\begin{solucion}{eje:2-jacobiano}
\datos $g(x,y)=(x^{2}y,\;3x-y,\;xy)\T$; punto $(1,2)$; $s=g_1+g_2+g_3=\uno\T g$.
\planteamiento cada fila de $J_g$ es el gradiente traspuesto de una componente; $\nabla s=J_g\T\uno$ (\cref{cor:cadena-escalar}).
\pasos
\begin{enumerate}[label=\arabic*.]
  \item Fila 1: $\nabla(x^{2}y)=(2xy,\,x^{2})=(4,1)$.
  \item Fila 2: $\nabla(3x-y)=(3,-1)$.
  \item Fila 3: $\nabla(xy)=(y,\,x)=(2,1)$.
  \item $J_g(1,2)=\begin{pmatrix}4&1\\3&-1\\2&1\end{pmatrix}$, de forma $3\times 2$.
  \item $\nabla s=J_g\T\uno=(4+3+2,\;1-1+1)\T=(9,1)\T$.
\end{enumerate}
\resultado \textbf{$J_g(1,2)$ es la matriz $3\times 2$ anterior y $\nabla s(1,2)=(9,1)\T$.}
\comprobacion $s=x^{2}y+3x-y+xy$: $s_x=2xy+3+y=4+3+2=9$ y $s_y=x^{2}-1+x=1$.
\end{solucion}

\begin{solucion}{eje:2-d2l}
\datos $w=f(x,y,z)$ con $x,y,z$ funciones de $(a,b)$; caso $f=xyz$, $x=a+b$, $y=a-b$, $z=ab$, en $(2,1)$.
\planteamiento \cref{cor:cadena-escalar} con tres variables intermedias.
\pasos
\begin{enumerate}[label=\arabic*.]
  \item (a) $\dfrac{\partial w}{\partial a}=\dfrac{\partial f}{\partial x}\dfrac{\partial x}{\partial a}+\dfrac{\partial f}{\partial y}\dfrac{\partial y}{\partial a}+\dfrac{\partial f}{\partial z}\dfrac{\partial z}{\partial a}$, y lo mismo con $b$.
  \item (b) Valores: $x=3$, $y=1$, $z=2$.
  \item Derivadas exteriores: $f_x=yz=2$, $f_y=xz=6$, $f_z=xy=3$.
  \item Respecto de $a$: $x_a=1$, $y_a=1$, $z_a=b=1$; luego $\partial w/\partial a=2+6+3=11$.
  \item Respecto de $b$: $x_b=1$, $y_b=-1$, $z_b=a=2$; luego $\partial w/\partial b=2-6+6=2$.
\end{enumerate}
\resultado \textbf{$\partial w/\partial a=11$ y $\partial w/\partial b=2$ en $(a,b)=(2,1)$.}
\comprobacion $w=(a^{2}-b^{2})ab=a^{3}b-ab^{3}$, de modo que $w_a=3a^{2}b-b^{3}=12-1=11$ y $w_b=a^{3}-3ab^{2}=8-6=2$.
\end{solucion}

\begin{solucion}{eje:2-coste}
\datos $J_1\in\R^{d\times n}$, $J_2\in\R^{d\times d}$, $J_3\in\R^{1\times d}$; $n=d=1000$.
\planteamiento un producto $(a\times b)(b\times c)$ cuesta $abc$ productos escalares (\cref{prop:modos}).
\pasos
\begin{enumerate}[label=\arabic*.]
  \item (a) $J_2J_1$: $d\cdot d\cdot n=d^{2}n$. Después $J_3(J_2J_1)$: $1\cdot d\cdot n=dn$. Total: $d^{2}n+dn$.
  \item (b) $J_3J_2$: $1\cdot d\cdot d=d^{2}$. Después $(J_3J_2)J_1$: $1\cdot d\cdot n=dn$. Total: $d^{2}+dn$.
  \item Con $n=d=1000$: (a) $10^{9}+10^{6}=1{,}001\cdot 10^{9}$; (b) $10^{6}+10^{6}=2\cdot 10^{6}$.
\end{enumerate}
\resultado \textbf{(a) $1{,}001\cdot 10^{9}$ productos (modo directo); (b) $2\cdot 10^{6}$ (modo inverso), unas 500 veces menos.}
\comprobacion $1{,}001\cdot 10^{9}/(2\cdot 10^{6})=500{,}5$. El modo inverso nunca forma una matriz $d\times n$ intermedia: solo multiplica vectores fila por matrices.
\end{solucion}

\begin{solucion}{eje:2-escalar}
\datos $a,b\colon\R^{n}\to\R^{m}$ diferenciables; $s=a\T b$; $q(x)=x\T Ax$.
\planteamiento escribir $s=\phi\circ G$ con $G(x)=(a(x),b(x))\in\R^{2m}$ y $\phi(\alpha,\beta)=\alpha\T\beta$, y aplicar el \cref{cor:cadena-escalar}.
\pasos
\begin{enumerate}[label=\arabic*.]
  \item $\phi(\alpha,\beta)=\sum_j\alpha_j\beta_j$, luego $\partial\phi/\partial\alpha_j=\beta_j$ y $\partial\phi/\partial\beta_j=\alpha_j$: $\nabla\phi=(\beta,\alpha)$.
  \item $J_G=\begin{pmatrix}J_a\\J_b\end{pmatrix}\in\R^{2m\times n}$, apilando los jacobianos por filas.
  \item $\nabla s=J_G\T\nabla\phi=\begin{pmatrix}J_a\T & J_b\T\end{pmatrix}\begin{pmatrix}b\\a\end{pmatrix}=J_a\T b+J_b\T a$.
  \item Forma cuadrática: $a(x)=x$ con $J_a=I$, y $b(x)=Ax$ con $J_b=A$ (\cref{ej:jac-lineal}).
  \item $\nabla q=I\,(Ax)+A\T x=(A+A\T)x$; si $A$ es simétrica, $\nabla q=2Ax$.
\end{enumerate}
\resultado \textbf{$\nabla(a\T b)=J_a\T b+J_b\T a$ y $\nabla(x\T Ax)=(A+A\T)x$, que vale $2Ax$ si $A=A\T$.} $\blacksquare$
\comprobacion con $A=\begin{pmatrix}1&2\\0&3\end{pmatrix}$, $q=x_1^{2}+2x_1x_2+3x_2^{2}$ y $\nabla q=(2x_1+2x_2,\,2x_1+6x_2)\T$; en $(1,1)$ vale $(4,8)\T$, y $(A+A\T)(1,1)\T=\begin{pmatrix}2&2\\2&6\end{pmatrix}(1,1)\T=(4,8)\T$.
\end{solucion}

% ---------------------------------------------------------------------
\section{Tema 3}

\begin{solucion}{eje:3-embeddings}
\datos $E\in\R^{50\,000\times 4096}$; $L$ escalar.
\planteamiento el gradiente respecto de una matriz tiene su misma forma (\cref{def:grad-matriz}).
\pasos
\begin{enumerate}[label=\arabic*.]
  \item $\partial L/\partial E$ tiene forma $50\,000\times 4096$.
  \item Número de entradas: $50\,000\cdot 4096=204\,800\,000$.
  \item La fila $i$ de $E$ solo interviene en $L$ si el token $i$ aparece en el lote; en caso contrario, $L$ no depende de ella y su gradiente es nulo.
\end{enumerate}
\resultado \textbf{forma $50\,000\times 4096$, con $2{,}048\cdot 10^{8}$ números; solo son no nulas las filas de los tokens presentes en el lote.}
\comprobacion orden de magnitud: $5\cdot 10^{4}\cdot 4\cdot 10^{3}=2\cdot 10^{8}$.
\end{solucion}

\begin{solucion}{eje:3-signos}
\datos $\nabla_zL=p-y$ con $y=e_c$.
\planteamiento estudiar el signo de cada componente y el efecto de restarla multiplicada por $\eta>0$.
\pasos
\begin{enumerate}[label=\arabic*.]
  \item Componente correcta: $p_c-1<0$, porque $0<p_c<1$ cuando $K\ge 2$.
  \item Componentes incorrectas: $p_i-0=p_i>0$.
  \item Paso de descenso: $z_c\leftarrow z_c+\eta(1-p_c)$ aumenta; $z_i\leftarrow z_i-\eta p_i$ disminuye para $i\ne c$, tanto más cuanto mayor es $p_i$.
\end{enumerate}
\resultado \textbf{la componente correcta es negativa y las demás positivas; el paso sube el logit correcto y baja los incorrectos en proporción a su probabilidad.}
\comprobacion en el \cref{ej:p-y}, $p-y=(-0{,}335;\,0{,}245;\,0{,}090)$: signos $(-,+,+)$ y suma nula.
\end{solucion}

\begin{solucion}{eje:3-capa}
\datos $W=\begin{pmatrix}1&0&-1\\2&1&0\end{pmatrix}$, $x=(1,2,1)\T$, $t=(0,3)\T$, $L=\tfrac12\norm{Wx-t}^{2}$.
\planteamiento $\delta=\nabla_y L=y-t$; después la \cref{prop:capa}.
\pasos
\begin{enumerate}[label=\arabic*.]
  \item $y_1=1-1=0$ e $y_2=2+2=4$: $y=(0,4)\T$.
  \item $\delta=y-t=(0,1)\T$ y $L=\tfrac12(0+1)=0{,}5$.
  \item $\partial L/\partial W=\delta x\T=\begin{pmatrix}0\\1\end{pmatrix}\begin{pmatrix}1&2&1\end{pmatrix}=\begin{pmatrix}0&0&0\\1&2&1\end{pmatrix}$.
  \item $\partial L/\partial x=W\T\delta=\begin{pmatrix}1&2\\0&1\\-1&0\end{pmatrix}\begin{pmatrix}0\\1\end{pmatrix}=(2,1,0)\T$.
\end{enumerate}
\resultado \textbf{$y=(0,4)\T$; $L=0{,}5$; $\partial L/\partial W=\begin{pmatrix}0&0&0\\1&2&1\end{pmatrix}$; $\partial L/\partial x=(2,1,0)\T$.}
\comprobacion las formas coinciden con las de $W$ ($2\times 3$) y $x$ (3 componentes). Como solo la segunda salida tiene error, $\partial L/\partial x$ es la segunda fila de $W$.
\end{solucion}

\begin{solucion}{eje:3-softmax}
\datos $z=(0,\ln 2,\ln 5)\T$; $c=2$, $y=(0,1,0)\T$.
\planteamiento $p_i=e^{z_i}/\sum_ke^{z_k}$, $L=-\ln p_c$ y \cref{teo:p-y}.
\pasos
\begin{enumerate}[label=\arabic*.]
  \item $e^{z}=(1,2,5)$, con suma $8$.
  \item $p=(1/8,\,1/4,\,5/8)=(0{,}125;\,0{,}25;\,0{,}625)$.
  \item $L=-\ln(1/4)=\ln 4=2\ln 2\approx 1{,}386$ nats.
  \item $\nabla_zL=p-y=(0{,}125;\,-0{,}75;\,0{,}625)$.
\end{enumerate}
\resultado \textbf{$p=(1/8,1/4,5/8)$; $L=\ln 4\approx 1{,}386$ nats; $\nabla_zL=(0{,}125;\,-0{,}75;\,0{,}625)$.}
\comprobacion $p$ suma~1 y $p-y$ suma $0{,}125-0{,}75+0{,}625=0$.
\end{solucion}

\begin{solucion}{eje:3-nucleo}
\datos $p=\softmax(z)$, $y=e_c$, $J=\diag(p)-pp\T$.
\planteamiento usar $\uno\T p=1$ y $\uno\T y=1$.
\pasos
\begin{enumerate}[label=\arabic*.]
  \item $\sum_j(p_j-y_j)=\uno\T p-\uno\T y=1-1=0$.
  \item $\diag(p)\uno=p$, porque la fila $i$ de $\diag(p)$ solo tiene la entrada $p_i$.
  \item $pp\T\uno=p\,(p\T\uno)=p$ por la asociatividad del producto.
  \item $J\uno=p-p=0$: $\uno\in\ker J$.
  \item $\nabla_zL\T\uno=\sum_j(p_j-y_j)=0$, luego el gradiente es ortogonal a $\uno$.
  \item Interpretación: $L(z+t\uno)=L(z)$ para todo $t$ (invariancia por traslación). Derivando en $t=0$ con el \cref{lem:curva}, $\nabla_zL\T\uno=0$, que es lo obtenido. Del mismo modo, $p(z+t\uno)$ es constante y su derivada, $J\uno$, es nula.
\end{enumerate}
\resultado \textbf{$\sum_j(p_j-y_j)=0$, $J\uno=0$ y $\nabla_zL\perp\uno$: la dirección $\uno$ no cambia ni las probabilidades ni la pérdida.} $\blacksquare$
\comprobacion para $z=(2,1,0)\T$, $J$ tiene autovalores $0$; $0{,}119$; $0{,}371$, y el autovector del autovalor nulo es $\uno/\sqrt3$.
\end{solucion}

\begin{solucion}{eje:3-word2vec}
\datos $s=U\T v_c$, con $s_w=u_w\T v_c$; $\hat y=\softmax(s)$; $J=-\ln\hat y_o$; datos numéricos del apartado (c).
\planteamiento por el \cref{teo:p-y}, $\partial J/\partial s=\hat y-y$; después, la regla de la cadena a través de $s=U\T v_c$, que es una capa lineal.
\pasos
\begin{enumerate}[label=\arabic*.]
  \item (a) $\partial s_w/\partial v_c=u_w$, luego $\partial J/\partial v_c=\sum_w(\hat y_w-y_w)\,u_w$.
  \item Una combinación lineal de las columnas de $U$ con coeficientes $\hat y-y$ es $U(\hat y-y)$. Es la \cref{prop:capa} con matriz $U\T$: $\partial J/\partial v_c=(U\T)\T\delta$.
  \item (b) $u_w$ solo aparece en $s_w$, con $\partial s_w/\partial u_w=v_c$; luego $\partial J/\partial u_w=(\hat y_w-y_w)\,v_c$.
  \item (c) $s_1=\ln 2$, $s_2=0$, $s_3=\ln 2$; $e^{s}=(2,1,2)$, con suma~5; $\hat y=(0{,}4;\,0{,}2;\,0{,}4)$.
  \item $J=-\ln 0{,}2=\ln 5\approx 1{,}609$ nats y $\hat y-y=(0{,}4;\,-0{,}8;\,0{,}4)$.
  \item $\partial J/\partial v_c=0{,}4\,(1,0)-0{,}8\,(0,1)+0{,}4\,(1,1)=(0{,}8;\,-0{,}4)$.
  \item Con $\ln 2\approx 0{,}6931$: $\partial J/\partial u_1=\partial J/\partial u_3=0{,}4\,(\ln 2,0)=(0{,}2773;\,0)$ y $\partial J/\partial u_2=-0{,}8\,(\ln 2,0)=(-0{,}5545;\,0)$.
\end{enumerate}
\resultado \textbf{$\partial J/\partial v_c=U(\hat y-y)=(0{,}8;\,-0{,}4)\T$; $\partial J/\partial u_w=(\hat y_w-y_w)v_c$: $(0{,}2773;\,0)$, $(-0{,}5545;\,0)$ y $(0{,}2773;\,0)$.} $\blacksquare$
\comprobacion la diferencia centrada con $h=10^{-6}$ reproduce estos valores con nueve cifras. Como la segunda componente de $v_c$ es nula, también lo es la de cada $\partial J/\partial u_w$.
\end{solucion}

% ---------------------------------------------------------------------
\section{Tema 4}

\begin{solucion}{eje:4-vf}
\planteamiento buscar contraejemplos para las falsas y aplicar el \cref{teo:schwarz} para la verdadera.
\pasos
\begin{enumerate}[label=(\alph*)]
  \item \textbf{Falsa.} $f(x,y)=x^{2}-y^{2}$ tiene $\nabla f(0,0)=0$, pero $f(t,0)=t^{2}>0$ y $f(0,t)=-t^{2}<0$: el origen es un punto de silla.
  \item \textbf{Verdadera}, por el \cref{teo:schwarz}.
  \item \textbf{Falsa.} $f(x,y)=-x^{4}+y^{2}$ tiene $\nabla f(0,0)=0$ y $H_f(0,0)=\diag(0,2)$, pero $f(t,0)=-t^{4}<0$ para $t\ne 0$.
\end{enumerate}
\resultado \textbf{(a) falsa; (b) verdadera; (c) falsa.}
\comprobacion en (c), con $x^{4}+y^{2}$ se obtiene el mismo hessiano en el origen y sí hay mínimo: el hessiano no decide (\cref{tab:criterio}).
\end{solucion}

\begin{solucion}{eje:4-taylor}
\datos $f(x,y)=x^{3}+xy$; punto $(1,1)$; $h=(0{,}1;\,-0{,}2)$.
\planteamiento \cref{teo:taylor}(ii) truncado en el primer y en el segundo orden; convexidad local mediante los autovalores de $H$.
\pasos
\begin{enumerate}[label=\arabic*.]
  \item $f(1,1)=2$ y $\nabla f=(3x^{2}+y,\,x)\T=(4,1)\T$.
  \item $\nabla f\T h=0{,}4-0{,}2=0{,}2$; primer orden: $2{,}2$.
  \item $H=\begin{pmatrix}6x&1\\1&0\end{pmatrix}=\begin{pmatrix}6&1\\1&0\end{pmatrix}$.
  \item $h\T Hh=6\cdot 0{,}01+2\cdot(0{,}1)(-0{,}2)=0{,}06-0{,}04=0{,}02$; segundo orden: $2{,}2+0{,}01=2{,}21$.
  \item Exacto: $f(1{,}1;\,0{,}8)=1{,}331+0{,}88=2{,}211$.
  \item $\det(H-\lambda I)=\lambda^{2}-6\lambda-1=0$ da $\lambda=3\pm\sqrt{10}$, es decir, $6{,}162$ y $-0{,}1623$.
\end{enumerate}
\resultado \textbf{primer orden $2{,}2$; segundo orden $2{,}21$; exacto $2{,}211$. $H(1,1)$ tiene un autovalor negativo, y por continuidad lo tiene en todo un entorno: $f$ no es convexa en ningún entorno de $(1,1)$.}
\comprobacion el error de segundo orden, $0{,}001=h_1^{3}$, es exactamente el término cúbico de $x^{3}$ que el desarrollo omite.
\end{solucion}

\begin{solucion}{eje:4-mit}
\datos $f(x,y)=x^{3}-3x^{2}+xy-y^{2}+x+y$.
\planteamiento resolver $\nabla f=0$ y aplicar el \cref{teo:criterio}.
\pasos
\begin{enumerate}[label=\arabic*.]
  \item $f_x=3x^{2}-6x+y+1$ y $f_y=x-2y+1$.
  \item De $f_y=0$: $y=(x+1)/2$. Sustituyendo en $f_x=0$ y multiplicando por~2: $6x^{2}-11x+3=0$.
  \item $x=(11\pm\sqrt{121-72})/12=(11\pm 7)/12$, es decir, $x=3/2$ o $x=1/3$; los puntos son $(3/2,\,5/4)$ y $(1/3,\,2/3)$.
  \item $H=\begin{pmatrix}6x-6&1\\1&-2\end{pmatrix}$.
  \item En $(3/2,5/4)$: $H=\begin{pmatrix}3&1\\1&-2\end{pmatrix}$, $\det H=-7<0$; autovalores $(1\pm\sqrt{29})/2\approx 3{,}193$ y $-2{,}193$: punto de silla.
  \item En $(1/3,2/3)$: $H=\begin{pmatrix}-4&1\\1&-2\end{pmatrix}$, $\det H=7>0$ y $f_{xx}=-4<0$; autovalores $-3\pm\sqrt2\approx-1{,}586$ y $-4{,}414$: máximo local.
  \item Valores: $f(3/2,5/4)=-5/16=-0{,}3125$ y $f(1/3,2/3)=13/27\approx 0{,}4815$.
\end{enumerate}
\resultado \textbf{$(3/2,5/4)$ es un punto de silla y $(1/3,2/3)$ un máximo local, con $f=13/27$.}
\comprobacion en $(1/3,2/3)$: $f_x=1/3-2+2/3+1=0$ y $f_y=1/3-4/3+1=0$.
\end{solucion}

\begin{solucion}{eje:4-tasa}
\datos $f(x,y)=x^{2}+4y^{2}=\tfrac12x\T Ax$ con $A=\diag(2,8)$; punto inicial $(1,1)$, con $f=5$.
\planteamiento \cref{prop:dg-cuadratica}: cada coordenada se multiplica por $1-\eta\lambda_i$ en cada paso.
\pasos
\begin{enumerate}[label=\arabic*.]
  \item (a) $\lambda_{\max}=8$: converge si y solo si $0<\eta<2/8=0{,}25$.
  \item (b) $\eta=0{,}2$: factores $1-0{,}4=0{,}6$ y $1-1{,}6=-0{,}6$. Iteraciones: $(0{,}6;\,-0{,}6)$ con $f=0{,}36+1{,}44=1{,}8$, y $(0{,}36;\,0{,}36)$ con $f=0{,}1296+0{,}5184=0{,}648$.
  \item $\eta=0{,}3$: factores $0{,}4$ y $1-2{,}4=-1{,}4$. Iteraciones: $(0{,}4;\,-1{,}4)$ con $f=0{,}16+7{,}84=8$, y $(0{,}16;\,1{,}96)$ con $f=0{,}0256+15{,}37=15{,}39$. La coordenada $y$ oscila y crece: diverge.
  \item (c) El peor factor $\max(\abs{1-2\eta},\abs{1-8\eta})$ es mínimo cuando $1-2\eta=-(1-8\eta)$, es decir, $\eta=2/10=0{,}2$, y vale $0{,}6$.
\end{enumerate}
\resultado \textbf{(a) $0<\eta<0{,}25$; (b) con $\eta=0{,}2$, $f$ pasa de $5$ a $1{,}8$ y a $0{,}648$; con $\eta=0{,}3$, a $8$ y a $15{,}39$; (c) $\eta=0{,}2$, con factor $0{,}6$.}
\comprobacion con $\kappa=8/2=4$, la fórmula de la observación sobre el número de condición da $(\kappa-1)/(\kappa+1)=3/5=0{,}6$; y en (b) cada paso multiplica $f$ por $0{,}6^{2}=0{,}36$: $5\cdot 0{,}36=1{,}8$.
\end{solucion}

\begin{solucion}{eje:4-necesaria}
\datos $f\in C^{2}(U)$; $x^{*}$ mínimo local.
\planteamiento por reducción al absurdo con el \cref{teo:taylor}(ii) a lo largo de un autovector.
\pasos
\begin{enumerate}[label=\arabic*.]
  \item Por el \cref{teo:fermat}, $\nabla f(x^{*})=0$.
  \item Supóngase que $H=H_f(x^{*})$ no es semidefinida positiva: existe un autovector unitario $v$ con autovalor $\lambda<0$.
  \item Por el \cref{teo:taylor}(ii), $f(x^{*}+tv)-f(x^{*})=\tfrac12\lambda t^{2}+o(t^{2})$.
  \item Para $\abs t$ pequeño y no nulo, el resto es menor que $\tfrac14\abs{\lambda}t^{2}$, luego $f(x^{*}+tv)-f(x^{*})\le\tfrac14\lambda t^{2}<0$, en contradicción con que $x^{*}$ es mínimo local.
  \item Contraejemplo del recíproco: $f(x,y)=x^{3}+y^{2}$ tiene $\nabla f(0,0)=0$ y $H_f(0,0)=\diag(0,2)\succeq 0$, pero $f(t,0)=t^{3}<0$ para $t<0$.
\end{enumerate}
\resultado \textbf{Si $x^{*}$ es mínimo local, $H_f(x^{*})\succeq 0$; el recíproco falla, por ejemplo, para $x^{3}+y^{2}$ en el origen.} $\blacksquare$
\comprobacion el contraejemplo es coherente con el \cref{teo:criterio}: el autovalor nulo deja el caso sin decidir.
\end{solucion}

% ---------------------------------------------------------------------
\section{Tema 5}

\begin{solucion}{eje:5-vf}
\planteamiento aplicar los teoremas del \cref{tema:5} o buscar un contraejemplo.
\pasos
\begin{enumerate}[label=(\alph*)]
  \item \textbf{Verdadera}: es el \cref{teo:segundo-orden} con $C=\R^{n}$, que es abierto y convexo.
  \item \textbf{Falsa.} $f(x)=x$ y $g(x)=x^{2}$ son convexas en $\R$, pero su producto $x^{3}$ tiene segunda derivada $6x<0$ para $x<0$.
  \item \textbf{Verdadera}, por la \cref{prop:operaciones}(iii).
\end{enumerate}
\resultado \textbf{(a) verdadera; (b) falsa; (c) verdadera.}
\comprobacion en (b), la cuerda de $x^{3}$ entre $-2$ y $0$ pasa por $(-1,-4)$, por debajo de la gráfica en $(-1,-1)$, lo que viola la \cref{def:convexa}.
\end{solucion}

\begin{solucion}{eje:5-hessianos}
\datos cuatro funciones de $\R^{2}$.
\planteamiento calcular el hessiano y sus autovalores (\cref{teo:segundo-orden}).
\pasos
\begin{enumerate}[label=(\alph*)]
  \item $H=\begin{pmatrix}2&1\\1&2\end{pmatrix}$, autovalores $1$ y $3$: estrictamente convexa.
  \item $H=\begin{pmatrix}2&-4\\-4&2\end{pmatrix}$, autovalores $2\pm 4$, es decir, $6$ y $-2$: no es convexa.
  \item $H=e^{x+y}\begin{pmatrix}1&1\\1&1\end{pmatrix}$, autovalores $0$ y $2e^{x+y}$: convexa. No es estrictamente convexa: es constante sobre cada recta $x+y=c$, donde la desigualdad de la \cref{def:convexa} es una igualdad.
  \item $H=\diag(12x^{2},2)\succeq 0$: convexa. Es estrictamente convexa, aunque $H$ sea singular en $x=0$: si $a\ne b$, o bien $a_1\ne b_1$ y la parte $x^{4}$ da desigualdad estricta, o bien $a_2\ne b_2$ y la da la parte $y^{2}$.
\end{enumerate}
\resultado \textbf{(a) estrictamente convexa; (b) no convexa; (c) convexa, no estrictamente; (d) estrictamente convexa.}
\comprobacion en (b), $f(t,t)=-2t^{2}$ es cóncava a lo largo de la diagonal, dirección del autovector $(1,1)$ de autovalor $-2$.
\end{solucion}

\begin{solucion}{eje:5-boyd}
\datos dominio $\R^{2}_{++}$, abierto y convexo.
\planteamiento \cref{teo:segundo-orden} aplicado a $f$ y a $-f$.
\pasos
\begin{enumerate}[label=(\alph*)]
  \item $f=x_1x_2$: $H=\begin{pmatrix}0&1\\1&0\end{pmatrix}$, autovalores $\pm1$. No es convexa ni cóncava.
  \item $f=x_1^{-1}x_2^{-1}$: $f_{11}=2x_1^{-3}x_2^{-1}$, $f_{12}=x_1^{-2}x_2^{-2}$, $f_{22}=2x_1^{-1}x_2^{-3}$. Entonces $f_{11}>0$ y $\det H=4x_1^{-4}x_2^{-4}-x_1^{-4}x_2^{-4}=3(x_1x_2)^{-4}>0$, luego $H\succ 0$: estrictamente convexa.
  \item $f=x_1^{2}/x_2$: $f_{11}=2/x_2$, $f_{12}=-2x_1/x_2^{2}$ y $f_{22}=2x_1^{2}/x_2^{3}$, es decir,
  \[H=\frac{2}{x_2^{3}}\begin{pmatrix}x_2^{2}&-x_1x_2\\-x_1x_2&x_1^{2}\end{pmatrix}=\frac{2}{x_2^{3}}\,ww\T,\qquad w=(x_2,-x_1)\T.\] Como $v\T ww\T v=(w\T v)^{2}\ge 0$ y $x_2>0$, $H\succeq 0$: convexa.
\end{enumerate}
\resultado \textbf{(a) ni convexa ni cóncava; (b) convexa (estrictamente); (c) convexa.}
\comprobacion en $(1,1)$: en (b), $H=\begin{pmatrix}2&1\\1&2\end{pmatrix}$, con autovalores $1$ y $3$; en (c), $H=\begin{pmatrix}2&-2\\-2&2\end{pmatrix}$, con autovalores $0$ y $4$.
\end{solucion}

\begin{solucion}{eje:5-red}
\datos $L(w_1,w_2)=(2-w_1w_2)^{2}$.
\planteamiento \cref{prop:simetria} y \cref{teo:criterio}.
\pasos
\begin{enumerate}[label=\arabic*.]
  \item (a) $L(1,2)=L(2,1)=(2-2)^{2}=0=\min L$, porque $L\ge 0$.
  \item Punto medio $(1{,}5;\,1{,}5)$: $L=(2-2{,}25)^{2}=0{,}0625>0$.
  \item (b) $L_{w_1}=-2w_2(2-w_1w_2)$, de donde $L_{w_1w_1}=2w_2^{2}$ y $L_{w_1w_2}=-2(2-w_1w_2)+2w_1w_2=4w_1w_2-4$; por simetría, $L_{w_2w_2}=2w_1^{2}$.
  \item $H_L(0,0)=\begin{pmatrix}0&-4\\-4&0\end{pmatrix}$, con autovalores $\pm 4$; como $\nabla L(0,0)=0$, el origen es un punto de silla.
  \item (c) El conjunto de mínimos no es convexo, luego $L$ no es convexa (\cref{prop:simetria}). Lo confirma también que $H_L(0,0)$ no es semidefinida positiva.
\end{enumerate}
\resultado \textbf{(a) $L(1,2)=L(2,1)=0$ y $L(1{,}5;\,1{,}5)=0{,}0625$; (b) silla, con autovalores $\pm4$; (c) $L$ no es convexa.}
\comprobacion la \cref{def:convexa} exigiría $L(1{,}5;\,1{,}5)\le\tfrac12\cdot 0+\tfrac12\cdot 0=0$, que es falso.
\end{solucion}

\begin{solucion}{eje:5-softmax-W}
\datos $L(W)=\varphi(Wx)$ con $\varphi(z)=-z_c+\LSE(z)$ convexa (\cref{cor:ce-convexa}) y $x$ fijo.
\planteamiento verificar la \cref{def:convexa} usando la linealidad de $W\mapsto Wx$.
\pasos
\begin{enumerate}[label=\arabic*.]
  \item Sean $W_1,W_2$, $\lambda\in[0,1]$, $z_1=W_1x$ y $z_2=W_2x$.
  \item Por la linealidad del producto, $\bigl(\lambda W_1+(1-\lambda)W_2\bigr)x=\lambda z_1+(1-\lambda)z_2$.
  \item Luego $L\bigl(\lambda W_1+(1-\lambda)W_2\bigr)=\varphi\bigl(\lambda z_1+(1-\lambda)z_2\bigr)\le\lambda\varphi(z_1)+(1-\lambda)\varphi(z_2)$, por la convexidad de $\varphi$.
  \item El último miembro es $\lambda L(W_1)+(1-\lambda)L(W_2)$, que es la desigualdad buscada.
  \item En una red $\varphi(W_2\,\sigma(W_1x))$, la aplicación $(W_1,W_2)\mapsto W_2\,\sigma(W_1x)$ no es lineal: aparece el producto de dos parámetros y la no linealidad $\sigma$. El paso~2 deja de valer y la desigualdad no puede encadenarse.
\end{enumerate}
\resultado \textbf{$L(W)$ es convexa en $W$; en una red con capa oculta falla el paso~2, y la convexidad se pierde (\cref{ej:red-lineal}).} $\blacksquare$
\comprobacion con $K=2$ y $m=z_c-z_{c'}$, $\varphi=\ln(1+e^{-m})$, cuya segunda derivada respecto de $m$ es $\sigma(m)(1-\sigma(m))\ge 0$, con $\sigma$ la función logística.
\end{solucion}

\begin{solucion}{eje:5-unico}
\datos $f$ estrictamente convexa en $C$; $L(z)=-z_c+\LSE(z)$; $L_\mu=L+\tfrac\mu2\norm{z}^{2}$.
\planteamiento reducción al absurdo con el punto medio, y \cref{prop:operaciones}.
\pasos
\begin{enumerate}[label=\arabic*.]
  \item Si $x_1\ne x_2$ fueran mínimos globales con valor $m$, la convexidad estricta daría $f\bigl(\tfrac12(x_1+x_2)\bigr)<\tfrac12m+\tfrac12m=m$, imposible.
  \item $L(z+t\uno)=L(z)$ para todo $t$; con $a=z$ y $b=z+\uno$, la desigualdad de la \cref{def:convexa} es una igualdad, luego $L$ no es estrictamente convexa.
  \item $\tfrac\mu2\norm{z}^{2}$ tiene hessiano $\mu I\succ 0$ y es estrictamente convexa; la suma de una función convexa y otra estrictamente convexa es estrictamente convexa, pues al sumar las dos desigualdades una de ellas es estricta.
  \item Por el paso~1, $L_\mu$ tiene a lo sumo un mínimo $z^{*}$. Si existe, $\nabla L_\mu(z^{*})=p-y+\mu z^{*}=0$ (\cref{teo:fermat}).
  \item Multiplicando por $\uno\T$: $\uno\T(p-y)+\mu\,\uno\T z^{*}=0+\mu\sum_iz^{*}_i=0$, luego $\sum_iz^{*}_i=0$.
\end{enumerate}
\resultado \textbf{Una función estrictamente convexa tiene a lo sumo un mínimo; $L$ no es estrictamente convexa por su invariancia en la dirección $\uno$, y $L_\mu$ sí lo es, con un mínimo (si existe) de componentes de suma nula.} $\blacksquare$
\comprobacion sin regularizar, $L$ ni siquiera alcanza su ínfimo: $L(t\,e_c)=-t+\ln(e^{t}+K-1)\to 0$ cuando $t\to\infty$, pero $L>0$ siempre.
\end{solucion}

% ---------------------------------------------------------------------
\section{Tema 6}

\begin{solucion}{eje:6-porque}
\datos un modelo con $N\sim 10^{10}$ parámetros.
\planteamiento comparar el coste de la \cref{prop:coste-fd} con el de la retropropagación (\cref{prop:modos}).
\pasos
\begin{enumerate}[label=\arabic*.]
  \item Las diferencias centradas requieren $2N$ evaluaciones de la pérdida por gradiente: unas $1{,}4\cdot 10^{10}$ para $N=7\cdot 10^{9}$.
  \item La retropropagación obtiene el gradiente exacto (salvo redondeo) con un coste de unas pocas evaluaciones, independiente de $N$.
  \item Además, las diferencias finitas tienen errores de truncamiento y de redondeo (\cref{prop:h-optimo}), y en bf16 son inútiles.
  \item Uso real: verificar la implementación de un gradiente en un modelo pequeño o en algunas componentes o direcciones aleatorias.
\end{enumerate}
\resultado \textbf{serían del orden de $10^{9}$ a $10^{10}$ veces más caras por paso y menos precisas; se usan solo para comprobar gradientes.}
\comprobacion con un coste de retropropagación de unas 3 evaluaciones, la razón es $1{,}4\cdot 10^{10}/3\approx 5\cdot 10^{9}$.
\end{solucion}

\begin{solucion}{eje:6-cubo}
\datos $f(x)=x^{3}$, $x=1$, $f'(1)=3$, $f''(1)=6$, $f'''=6$.
\planteamiento \cref{def:diferencias}; los errores previstos son $\tfrac h2f''(\xi)$ y $\tfrac{h^{2}}{6}f'''=h^{2}$.
\pasos
\begin{enumerate}[label=\arabic*.]
  \item $h=0{,}1$: $f(1{,}1)=1{,}331$ y $f(0{,}9)=0{,}729$.
  \item Hacia delante: $0{,}331/0{,}1=3{,}31$, error $0{,}31$. Centrada: $0{,}602/0{,}2=3{,}01$, error $0{,}01$.
  \item $h=0{,}01$: $f(1{,}01)=1{,}030301$ y $f(0{,}99)=0{,}970299$ (se conservan todas las cifras, porque la resta las necesita).
  \item Hacia delante: $0{,}030301/0{,}01=3{,}0301$, error $0{,}0301$. Centrada: $0{,}060002/0{,}02=3{,}0001$, error $0{,}0001$.
\end{enumerate}
\resultado \textbf{al dividir $h$ entre~10, el error hacia delante se divide aproximadamente entre~10 ($0{,}31\to 0{,}0301$) y el centrado exactamente entre~100 ($0{,}01\to 0{,}0001$).}
\comprobacion para un cúbico, $f'''$ es constante y la \cref{prop:truncamiento}(ii) da el error exacto $h^{2}$: $0{,}1^{2}=0{,}01$ y $0{,}01^{2}=0{,}0001$. El error hacia delante es $3h+h^{2}$: $0{,}31$ y $0{,}0301$.
\end{solucion}

\begin{solucion}{eje:6-direccional}
\datos $f(x,y)=x^{2}y+3y$; $x=(2,1)$; $v=(1,1)\T$; $h=0{,}01$.
\planteamiento \cref{prop:coste-fd}: comparar $\nabla f\T v$ con $\bigl(f(x+hv)-f(x-hv)\bigr)/(2h)$.
\pasos
\begin{enumerate}[label=\arabic*.]
  \item Analítico: $\nabla f(2,1)\T v=4+7=11$.
  \item $f(2{,}01;\,1{,}01)=4{,}0401\cdot 1{,}01+3{,}03=4{,}080501+3{,}03=7{,}110501$.
  \item $f(1{,}99;\,0{,}99)=3{,}9601\cdot 0{,}99+2{,}97=3{,}920499+2{,}97=6{,}890499$.
  \item Numérico: $(7{,}110501-6{,}890499)/0{,}02=0{,}220002/0{,}02=11{,}0001$.
  \item Error relativo: $0{,}0001/11{,}0001\approx 9{,}1\cdot 10^{-6}$.
\end{enumerate}
\resultado \textbf{$11{,}0001$ frente a $11$, con error relativo $9{,}1\cdot 10^{-6}$: el gradiente es correcto; la discrepancia es de truncamiento.}
\comprobacion $\phi(t)=f(2+t,1+t)=t^{3}+5t^{2}+8t+4+3(1+t)$ tiene $\phi'''=6$, y el error previsto es $\tfrac{h^{2}}{6}\cdot 6=h^{2}=10^{-4}$, exactamente el observado. Con $h=10^{-5}$ el error relativo bajaría a unos $10^{-11}$.
\end{solucion}

\begin{solucion}{eje:6-cs231n}
\datos tres pares (analítico, numérico) y $f(x)=\max(0,x)$, $x=10^{-6}$, $h=10^{-5}$.
\planteamiento \cref{def:error-relativo} y criterios orientativos de~\cite{cs231n}.
\pasos
\begin{enumerate}[label=(\alph*)]
  \item $r=10^{-7}/1{,}0000001\approx 1{,}0\cdot 10^{-7}$: gradiente correcto.
  \item $r=0{,}02/0{,}52\approx 3{,}8\cdot 10^{-2}>10^{-2}$: casi con seguridad hay un error.
  \item $r=5\cdot 10^{-7}/1{,}0005\cdot 10^{-3}\approx 5{,}0\cdot 10^{-4}$: zona dudosa; en una función derivable debe revisarse. Obsérvese que el error absoluto, $5\cdot 10^{-7}$, parecería pequeño: por eso se usa el relativo.
  \item $f(x+h)=1{,}1\cdot 10^{-5}$ y $f(x-h)=\max(0,-9\cdot 10^{-6})=0$; la diferencia centrada es $1{,}1\cdot 10^{-5}/(2\cdot 10^{-5})=0{,}55$, frente a $f'(x)=1$. Error relativo: $0{,}45/1=0{,}45$.
\end{enumerate}
\resultado \textbf{(a) $1{,}0\cdot 10^{-7}$, correcto; (b) $3{,}8\cdot 10^{-2}$, erróneo; (c) $5{,}0\cdot 10^{-4}$, sospechoso; (d) $0{,}55$ frente a~1, error relativo $0{,}45$, porque $x-h<0<x+h$: el intervalo cruza el punto no derivable de ReLU.}
\comprobacion en (d), con $x=10^{-3}$ el intervalo $[x-h,x+h]$ queda en la parte lineal y la diferencia centrada vale exactamente~1.
\end{solucion}

\begin{solucion}{eje:6-fermi}
\datos $N=7\cdot 10^{9}$; 1~s por evaluación.
\planteamiento \cref{prop:coste-fd}: $2N$ evaluaciones.
\pasos
\begin{enumerate}[label=\arabic*.]
  \item Evaluaciones: $2\cdot 7\cdot 10^{9}=1{,}4\cdot 10^{10}$.
  \item Un año tiene $365{,}25\cdot 24\cdot 3600=3{,}156\cdot 10^{7}$~s.
  \item Tiempo: $1{,}4\cdot 10^{10}/3{,}156\cdot 10^{7}\approx 443{,}6$ años, para un solo gradiente.
  \item Alternativa: retropropagación para el gradiente y comprobación numérica en un modelo diminuto o en unas pocas direcciones aleatorias, a dos evaluaciones por dirección.
\end{enumerate}
\resultado \textbf{$1{,}4\cdot 10^{10}$ evaluaciones, unos 444 años por gradiente.}
\comprobacion $10^{10}\text{ s}/(3\cdot 10^{7}\text{ s/año})\approx 330$ años, que multiplicado por $1{,}4$ da del orden de 450.
\end{solucion}

\begin{solucion}{eje:6-hoptimo}
\datos errores de evaluación $\le\varepsilon\abs f$; $\abs{\phi'''}\le M$, $\abs{\phi''}\le M_2$; $f=e^{x}$ en $0$ con $\abs f=M=1$, $\varepsilon=1{,}1\cdot 10^{-16}$.
\planteamiento acotar por separado redondeo y truncamiento, y minimizar la suma en $h>0$.
\pasos
\begin{enumerate}[label=\arabic*.]
  \item Redondeo (centrada): el numerador calculado es $f(x+h)-f(x-h)+\epsilon_1-\epsilon_2$ con $\abs{\epsilon_i}\le\varepsilon\abs f$; dividido por $2h$, el error es $\le 2\varepsilon\abs f/(2h)=\varepsilon\abs f/h$.
  \item Truncamiento: $\le Mh^{2}/6$ por la \cref{prop:truncamiento}(ii). Total: $E(h)=\varepsilon\abs f/h+Mh^{2}/6$.
  \item $E'(h)=-\varepsilon\abs f/h^{2}+Mh/3=0$ si y solo si $h^{3}=3\varepsilon\abs f/M$; $E'$ es negativa antes y positiva después, luego $h^{*}=(3\varepsilon\abs f/M)^{1/3}$ es el mínimo global.
  \item Con $\varepsilon\abs f=M(h^{*})^{3}/3$: $E(h^{*})=M(h^{*})^{2}/3+M(h^{*})^{2}/6=M(h^{*})^{2}/2$.
  \item Caso numérico: $h^{*}=(3{,}3\cdot 10^{-16})^{1/3}\approx 6{,}91\cdot 10^{-6}$ y $E(h^{*})=\tfrac12(6{,}91\cdot 10^{-6})^{2}\approx 2{,}39\cdot 10^{-11}$.
  \item Hacia delante: $E_+(h)=2\varepsilon\abs f/h+M_2h/2$; $E_+'(h)=0$ da $h_+^{*}=2(\varepsilon\abs f/M_2)^{1/2}$ y $E_+(h_+^{*})=2(\varepsilon\abs fM_2)^{1/2}$, ambos de orden $\varepsilon^{1/2}$. Con los datos, $h_+^{*}\approx 2{,}1\cdot 10^{-8}$ y $E_+\approx 2{,}1\cdot 10^{-8}$.
\end{enumerate}
\resultado \textbf{$h^{*}=(3\varepsilon\abs f/M)^{1/3}\approx 6{,}9\cdot 10^{-6}$ y $E(h^{*})=\tfrac M2(h^{*})^{2}\approx 2{,}4\cdot 10^{-11}$; hacia delante, paso y error de orden $\varepsilon^{1/2}\approx 10^{-8}$.} $\blacksquare$
\comprobacion en la \cref{fig:error-h}, el error medido mínimo es $1{,}2\cdot 10^{-11}$ en $h=10^{-5}$, por debajo de la cota $2{,}4\cdot 10^{-11}$ y con un paso próximo a $h^{*}$. La diferencia centrada es unas mil veces más precisa que la diferencia hacia delante en su mejor caso.
\end{solucion}
